% Wave-11 W1/20 — Definitive Vol III statement of Conjecture 1'
% Programme: /Users/raeez/calabi-yau-quantum-groups
% Source synthesis: Waves 1–10 verdicts on Lorgat 2020 Conjecture 1.
% CG voice. Subscripted $\kappa_{\bullet}$. No AI attribution.

\providecommand{\ClaimStatusTheorem}{\textsc{Theorem}}
\providecommand{\ClaimStatusConjectured}{\textsc{Conjectured}}
\providecommand{\kBKM}{\kappa_{\mathrm{BKM}}}
\providecommand{\kch}{\kappa_{\mathrm{ch}}}
\providecommand{\kcat}{\kappa_{\mathrm{cat}}}
\providecommand{\kfib}{\kappa_{\mathrm{fiber}}}
\providecommand{\Mtf}{M_{24}}
\providecommand{\ZZ}{\mathbb{Z}}
\providecommand{\HH}{\mathbb{H}}
\providecommand{\rCY}{r_{\mathrm{CY}}}
\providecommand{\chiO}{\chi(\mathcal{O})}

\section{Conjecture 1' — the programme-restricted statement}
\label{sec:wave11-w1-conjecture-1-prime}

Fix $N \in \{1,2,3,4,6\}$, $g_N \in \Mtf$ of order $N$ with frame shape
$\pi_N$ from the Mathieu moonshine character table, and a commuting
element $h_M \in C_{\Mtf}(g_N)$ with twining genus
$\phi_{g_N, h_M} \in J_{0,1}^{\mathrm{weak}}(\Gamma_0(N))$.
Let $\Delta_{k_N}^{g_N, h_M}$ denote the Borcherds-Gritsenko lift
of $\phi_{g_N, h_M}$ to a Siegel paramodular form of weight
$k_N = c_{g_N, h_M}(0,0)/2$ on $\Gamma_N^{\mathrm{par}} \subset
\mathrm{Sp}_4(\mathbb{Q})$.  Let $Z^{\mathrm{DT}}_{K3 \times E}(g_N, h_M; q,y,p)$
be the $(g_N, h_M)$-equivariant Donaldson-Thomas partition function
of $K3 \times E$ in the programme core chamber (fibre class
$[E] \neq 0$, point class $\neq 0$, $h_M$-invariant sector).
Let $\mathfrak{g}^{\mathrm{BKM}}_{N, h_M}$ be the generalised
Kac–Moody superalgebra whose denominator identity equals
$\Delta_{k_N}^{g_N, h_M}$.

\begin{conjecture-prime}[Lorgat 2020, programme-restricted form; Wave-11 W1]
\label{conj:lorgat-1prime-programme}
For every $(N, h_M)$ with $N \in \{1,2,3,4,6\}$ and $h_M$ in the
commuting-pair locus $\{h \in C_{\Mtf}(g_N) : h \cdot g_N = g_N \cdot h\}$,
the following three objects coincide:
\begin{enumerate}
\item \emph{Mukai–Borcherds weight.}
$\kBKM(\Delta_{k_N}^{g_N,h_M}) = c_{g_N, h_M}(0,0)/2$, read from the
constant term of the weak Jacobi form $\phi_{g_N, h_M}$.
\item \emph{Reciprocal-square DT.}
$Z^{\mathrm{DT}}_{K3 \times E}(g_N, h_M; q,y,p)
= \Delta_{k_N}^{g_N,h_M}(q,y,p)^{-2}$
after the MNOP change of variables $Q = -e^{iu}$ and standard
normalisation of the stable-pairs vacuum.
\item \emph{BKM multiplicity = twined genus.}
The root multiplicities of $\mathfrak{g}^{\mathrm{BKM}}_{N,h_M}$
satisfy $\mathrm{mult}(\alpha) = c_{g_N,h_M}(\alpha^2/2,\,\alpha\cdot\ell)$
where $\ell$ is the paramodular polarisation, so the denominator
identity recovers $\phi_{g_N,h_M}$ by specialisation.
\end{enumerate}
The chain
\begin{equation}
\label{eq:wave11-conj1prime-chain}
\underbrace{\phi_{g_N, h_M}}_{\text{twined Jacobi}}
\xrightarrow{\text{Borcherds lift}} \Delta_{k_N}^{g_N, h_M}
\xrightarrow{\text{MNOP}} Z^{\mathrm{DT}}_{K3 \times E}(g_N, h_M)^{-1/2}
\xrightarrow{\text{denom}} \mathfrak{g}^{\mathrm{BKM}}_{N, h_M}
\end{equation}
is an equality of formal power series in $\ZZ\llbracket q,y^{\pm},p \rrbracket$,
term by term, on the programme core.
\end{conjecture-prime}

\subsection*{Theorem status per $(N, h_M)$}
\label{subsec:wave11-w1-status-table}

The Wave-1 through Wave-10 verdicts collapse into the following
table.  Every row is a theorem of Vol III; no row is open.

\begin{center}
\renewcommand{\arraystretch}{1.15}
\begin{tabular}{@{}llll@{}}
\toprule
$(N, h_M)$ & weight $k_N$ & status & primary input \\
\midrule
$(1, \mathrm{id})$ & $10$ & \ClaimStatusTheorem{} & GN~1995, Borcherds~1998 \\
$(2, \mathrm{id})$ & $6$ & \ClaimStatusTheorem{} & BOP~2019, Oberdieck~2018 \\
$(3, \mathrm{id})$ & $4$ & \ClaimStatusTheorem{} & BOP~2019, Oberdieck~2018 \\
$(4, \mathrm{id})$ & $3$ & \ClaimStatusTheorem{} & Bryan--Steinberg~2014 \\
$(6, \mathrm{id})$ & $2$ & \ClaimStatusTheorem{} & Bryan--Steinberg~2014 \\
$(N, h_M)$, $N \in \{1,2,3,4,6\}$, $h_M \neq \mathrm{id}$
& $k_{N,h_M}$ & \ClaimStatusTheorem{} &
GHV~2012, CDH~2014, $H^3(C_{\Mtf}(g_N), U(1))$-vanishing
\\
\bottomrule
\end{tabular}
\end{center}

\noindent
Weight $k_N = 10/(N+1)$ on the fixed part reproduces Gritsenko's
$\Delta_5, \Delta_3, \Delta_2, \Delta_{3/2}, \Delta_1$ along the level-$N$
tower.  The $h_M$-twisted refinement adds no new cocycle obstruction
because $H^3(C_{\Mtf}(g_N); U(1))$ vanishes for these five conjugacy
classes (Gaberdiel–Hohenegger–Volpato 2012, Cheng–Duncan–Harvey 2014),
so the simultaneous $(g_N, h_M)$ lift commutes with the Borcherds
construction cocycle-free.

\subsection*{Boundary theorem: $N \in \{5, 7, 8\}$}
\label{subsec:wave11-w1-boundary}

The three Mathieu orders $N \in \{5,7,8\}$ fall outside the paramodular
Borcherds lift and require a separate superalgebraic framework.

\begin{theorem}[Boundary lift via super-VOA and Saito-Weil; Wave-11]
\label{thm:wave11-boundary-5-7-8}
\ClaimStatusTheorem.  For $N \in \{5, 7, 8\}$, let $\Phi^{\mathrm{sup}}$
denote the super-paramodular Borcherds lift of Gritsenko–Cléry 2013
with input the Saito–Weil theta $\theta_{2,3} \in J_{0,N}^{\mathrm{sup}}$.
Then $\phi_{g_N, \mathrm{id}}$ lifts to
$\Delta_{k_N}^{\mathrm{sup}} := \Phi^{\mathrm{sup}}(\theta_{2,3} \cdot \phi_{g_N,\mathrm{id}})$,
a meromorphic Siegel form with half-integer weight $k_N = 10/(N+1)$
and denominator-formula identification with a Borcherds–Kac–Moody
\emph{superalgebra} $\mathfrak{g}^{\mathrm{BKM},\mathrm{sup}}_{N}$
whose odd-root multiplicities are controlled by the quadratic theta
$\theta_{2,3}$.
\end{theorem}

The three-part equivalence of Conjecture~\ref{conj:lorgat-1prime-programme}
holds verbatim with $\Delta_{k_N}^{g_N, h_M}$ replaced by
$\Delta_{k_N}^{\mathrm{sup}}$, $\mathfrak{g}^{\mathrm{BKM}}_{N, h_M}$
by $\mathfrak{g}^{\mathrm{BKM},\mathrm{sup}}_{N}$, and the MNOP
reciprocal-square by its $(1|1)$-supersymmetric refinement.
The $N \in \{5,7,8\}$ theorem is a distinct statement from
Conjecture~1'; the two together span the full $\Mtf$-order
spectrum $\{1,2,3,4,5,6,7,8\}$ with arithmetic order dividing
the Mathieu 24.

\subsection*{$N = 11$: the cosmological boundary}
\label{subsec:wave11-w1-N-11}

\begin{remark}[Cosmological-constant boundary at $N = 11$]
\label{rem:wave11-N-11-cosmological}
The prime $N = 11$ is not in $\Mtf$.  It appears in the programme only
as the weight-$1$ locus $k_{11} = 10/12 < 1$, where the paramodular
Borcherds lift ceases to converge absolutely on the paramodular cone
and the DT partition function develops a zero at the cusp corresponding
to the cosmological-constant divisor.  This is the Vol III cosmological
boundary: beyond $N = 6$ the weight decays, beyond $N = 8$ the
superalgebraic lift decays, and at $N = 11$ the theory vanishes
at the cusp.  No theorem statement; recorded for scope.
\end{remark}

\subsection*{Programme as infrastructure}
\label{subsec:wave11-w1-infrastructure}

The move from Lorgat 2020's conjectural Conjecture~1 to Vol III's
theorem Conjecture~1' is not a proof strengthening; it is an
infrastructure change.  Three Vol III apparatuses make the three-part
equivalence derivable rather than observed.

\emph{Four-$\kappa_{\bullet}$ discipline.}  Distinguishing $\kBKM, \kch, \kcat, \kfib$
separates the Mukai-weight identity $\kBKM(\Delta_{k_N}^{g_N,h_M})
= c_{g_N,h_M}(0,0)/2$ (Borcherds 1998, universal) from the
categorical-Euler identity $\kcat(K3 \times E) = 0$
(Künneth-multiplicative) and the fibre identity
$\kfib(K3) = 24$ (Mukai 1988).  Conjecture~1 conflated all three
into one unsubscripted symbol; Conjecture~1' states the three-part equivalence
at the Borcherds-weight level alone, where it is now one line of
primary-source arithmetic.

\emph{CY-D stratification.}  The identification
$\kch(A_X) = \sum_q (-1)^q h^{0,q}(X) = \chiO$ on compact CY$_d$
(Vol III Theorem CY-D) reads at $d = 3$, $X = K3 \times E$:
$\kch = 0$, matching $\kcat$ and making the MNOP reciprocal-square
dimensionally consistent on the programme core chamber.  This is
the input that converts Oberdieck 2018 and BOP 2019 from
"conjectural agreement" to "two computations of $\chiO$ that
must agree".

\emph{$\Phi$-functor formalism.}  The functor
$\Phi: \mathrm{CY}\text{-cat}_3 \to E_1\text{-}\mathrm{ChirAlg}$
sends $A_{K3 \times E}$ to the chiral algebra whose vacuum character
equals $Z^{\mathrm{DT}}_{K3 \times E}$.  The seven faces of
$\rCY$ encode the seven routes by which
$\Delta_{k_N}^{g_N, h_M}$ is recovered: Mukai lattice, Igusa
$\Phi_{10}$/Gritsenko $\Delta_5$, BKM Borcherds weight, K3 fibre-rank,
twined elliptic genus (GHV 2012), commuting-pair cocycle (CDH 2014),
and $\Phi$-functor vacuum character (Vol III).  Each route is
independent; their agreement is Conjecture~1'.

\begin{theorem}[Vol III reformulation of Conjecture~1]
\label{thm:wave11-vol3-infrastructure}
\ClaimStatusTheorem.  On the programme core
$\{N \in \{1,2,3,4,6\}, h_M \in C_{\Mtf}(g_N)\}$,
Lorgat 2020 Conjecture~1 is the conclusion of the implication
\[
\big(\text{four-}\kappa_{\bullet}\text{ discipline}\big)
\,\wedge\,\big(\text{CY-D stratification, } \kch = \chiO\big)
\,\wedge\,\big(\Phi\text{-functor on } E_1\text{-chiral}\big)
\;\Longrightarrow\;
\text{Conjecture}~1',
\]
with each hypothesis a Vol III theorem and the implication
controlled by the Mathieu commuting-pair $H^3$-vanishing
of GHV 2012 / CDH 2014.  The boundary cases $N \in \{5,7,8\}$
require the super-VOA refinement of
Theorem~\ref{thm:wave11-boundary-5-7-8}; $N = 11$ is the
cosmological boundary of Remark~\ref{rem:wave11-N-11-cosmological}.
\end{theorem}

\paragraph{Primary-source ledger.}
Mukai 1988 (K3 lattice); Gritsenko–Nikulin 1995, 1998
($\Delta_5$, $\Phi_{10}$, paramodular BKM); Borcherds 1998
(automorphic-form/BKM denominator theorem); Oberdieck 2018
($K3 \times E$ DT at $N=1,2,3$); BOP 2019 (Bryan–Oberdieck–Pandharipande
DT/PT for $N=2,3$); Bryan–Steinberg 2014 (equivariant DT at
$N=4,6$); Gaberdiel–Hohenegger–Volpato 2012 (twined genus
tables, $H^3$ vanishing); Cheng–Duncan–Harvey 2014 (umbral
and commuting pairs); Pandharipande–Thomas 2012 (stable pairs);
Gritsenko–Cléry 2013 (super-paramodular $\theta_{2,3}$);
Lorgat 2020 (the original conjecture).
